% ============================================================
% STATUT D'AUDIT — ce fichier est dans sa version D'ORIGINE.
% Les reparations demandees par les audits n'y sont PAS appliquees.
% Voir maths/JOURNAL.md, section correspondante, avant tout usage.
% ============================================================
% =====================================================================
% Publication-grade Gaussian-corner asymptotic and B1 horizontal coupling
% =====================================================================

\subsection{A uniform Gaussian-corner theorem}

Write
\[
 \varphi(t)=(2\pi)^{-1/2}e^{-t^2/2},
 \qquad
 r_z=-\Phi^{-1}(z),\qquad 0<z<1/2.
\]
The next elementary lemma records the normal-tail estimates used below.

\begin{lemma}[Uniform normal-tail coordinates]
\label{lem:uniform-normal-coordinate}
For \(q>0\), put
\[
 \eta(q)=\log\left\{\frac{q\Phi(-q)}{\varphi(q)}\right\}.
\]
Then
\begin{equation}
 -\log(1+q^{-2})\le \eta(q)\le0.
 \label{eq:normal-mills-log}
\end{equation}
Consequently,
\begin{equation}
 \log\Phi(-q)
 =
 -\frac{q^2}{2}-\log(\sqrt{2\pi}q)+\eta(q).
 \label{eq:normal-log-tail}
\end{equation}

Fix \(M<\infty\). For \(r>0\), \(z=\Phi(-r)\), and
\(0<x\le M\), define
\[
 q_r(x)=-\Phi^{-1}(zx).
\]
There are constants \(r_0(M)\) and \(C_M<\infty\) such that, for
\(r\ge r_0(M)\),

\begin{enumerate}
\item
\begin{equation}
 q_r(x)\ge r-\frac{C_M}{r},
 \qquad 0<x\le M;
 \label{eq:q-lower-uniform}
\end{equation}

\item whenever \(|\log x|\le\sqrt r\),
\begin{equation}
 q_r(x)
 =
 r-\frac{\log x}{r}+\epsilon_r(x),
 \qquad
 |\epsilon_r(x)|
 \le
 \frac{C_M\{1+|\log x|+(\log x)^2\}}{r^3}.
 \label{eq:q-expansion-uniform}
\end{equation}
\end{enumerate}
All constants are deterministic.
\end{lemma}

\begin{proof}
The classical Mills inequalities
\[
 \frac{q}{1+q^2}\varphi(q)
 \le \Phi(-q)
 \le \frac{\varphi(q)}q
\]
give \eqref{eq:normal-mills-log} and
\eqref{eq:normal-log-tail}.

Let \(\ell=\log x\), \(q=q_r(x)\). Since
\(\Phi(-q)=x\Phi(-r)\), \eqref{eq:normal-log-tail} gives
\begin{equation}
 \ell
 =
 -\frac{q^2-r^2}{2}
 -\log\frac qr
 +\eta(q)-\eta(r).
 \label{eq:q-implicit}
\end{equation}
For fixed \(M\), Mills' inequalities imply
\(\Phi(-r/2)/\Phi(-r)\to\infty\), so \(q_r(M)\ge r/2\) for
large \(r\). By monotonicity, \(q_r(x)\ge q_r(M)\ge r/2\) for
all \(x\le M\).

Suppose now that \(|\ell|\le\sqrt r\), and put
\[
 q_0=r-\frac{\ell}{r}.
\]
Then
\begin{align}
 -\frac{q_0^2-r^2}{2}-\log\frac{q_0}{r}
 &=
 \ell-\frac{\ell^2}{2r^2}
 -\log\left(1-\frac{\ell}{r^2}\right)
 \nonumber\\
 &=
 \ell
 +O\left(\frac{|\ell|+\ell^2}{r^2}\right).
 \label{eq:q0-residual}
\end{align}
Moreover \(q_0\ge r/2\) for large \(r\), and
\eqref{eq:normal-mills-log} gives
\[
 |\eta(q_0)-\eta(r)|\le \frac{C}{r^2}.
\]
Thus
\[
 \left|
 \log\frac{\Phi(-q_0)}{\Phi(-r)}-\ell
 \right|
 \le
 \frac{C_M\{1+|\ell|+\ell^2\}}{r^2}.
\]
The derivative of \(s\mapsto\log\Phi(-s)\) is
\[
 -\frac{\varphi(s)}{\Phi(-s)},
\]
whose absolute value is at least \(s\) by the upper Mills
inequality. The entire interval between \(q\) and \(q_0\) lies
above \(r/3\) for large \(r\). The mean-value theorem therefore
gives
\[
 |q-q_0|
 \le
 \frac{C_M\{1+|\ell|+\ell^2\}}{r^3},
\]
which is \eqref{eq:q-expansion-uniform}.

Taking \(x=M\) in \eqref{eq:q-expansion-uniform}, and using
monotonicity \(q_r(x)\ge q_r(M)\), proves
\eqref{eq:q-lower-uniform}.
\end{proof}


\begin{theorem}[Uniform Gaussian-corner asymptotic]
\label{thm:uniform-gaussian-corner}
Let \(1\le d\le d_0\), and let
\[
 X\sim N_d(\mu,\Sigma),
 \qquad
 \Omega=\Sigma^{-1},
 \qquad
 U=\Phi(X)
\]
with \(\Phi\) applied coordinatewise. Let the tuple
\[
 (d,\mu,\Sigma,D,a,w)
\]
range over a family satisfying the following uniform conditions.

\begin{enumerate}
\item[(G1)]
There are constants
\[
 0<\omega_-\le\omega_+<\infty,
 \qquad M_\mu<\infty,
 \qquad \nu_0>0,
\]
such that
\[
 \omega_-\le\lambda_{\min}(\Omega)
 \le\lambda_{\max}(\Omega)\le\omega_+,
 \qquad
 \|\mu\|_2\le M_\mu,
\]
and, with
\[
 \nu=\Omega\mathbf 1,
 \qquad
 \kappa=\mathbf 1^\top\Omega\mathbf 1,
 \qquad
 b=\mathbf 1^\top\Omega\mu,
\]
one has
\[
 \min_{1\le i\le d}\nu_i\ge\nu_0.
\]

\item[(G2)]
There are constants
\[
 \eta_0>0,\quad
 M_D<\infty,\quad
 0<a_-\le a_+<\infty,\quad
 0<c_*\le C_*<\infty,\quad
 d_*>0
\]
such that \(D:[0,1]^d\to[0,\infty)\) satisfies
\[
 D(0)=0,\qquad
 \nabla D(0)=a=(a_1,\ldots,a_d),
 \qquad
 a_-\le a_i\le a_+,
\]
\[
 \sup_{u\in[0,\eta_0]^d}
 \|\nabla^2D(u)\|_{\mathrm{op}}\le M_D,
\]
\[
 c_*\|u\|_1\le D(u)\le C_*\|u\|_1,
 \qquad u\in[0,\eta_0]^d,
\]
and
\[
 \inf_{u\notin[0,\eta_0]^d}D(u)\ge d_*.
\]

\item[(G3)]
The weight \(w\) is positive and measurable, and there are
\(0<w_-\le w_+<\infty\) and \(L_w<\infty\) such that
\[
 w_-\le w(0)\le w_+,
 \qquad
 |w(u)-w(0)|\le L_w\|u\|_1,
 \qquad u\in[0,\eta_0]^d.
\]
The unweighted case corresponds to \(w\equiv1\).
\end{enumerate}

Then, uniformly over this family, as \(z\downarrow0\),
\begin{align}
 &\E\left[
   w\{\Phi(X)\}
   \mathbf 1_{\{D(\Phi(X))\le z\}}
 \right]
 \nonumber\\
 &\quad=
 w(0)\,C(\mu,\Sigma,a)\,
 z^\kappa
 \{\sqrt{2\pi}r_z\}^{\kappa-d}
 e^{-br_z}
 \left\{1+O(r_z^{-1})\right\},
 \label{eq:uniform-corner-main}
\end{align}
where
\begin{equation}
 C(\mu,\Sigma,a)
 =
 |\Sigma|^{-1/2}
 e^{-\mu^\top\Omega\mu/2}
 \frac{\prod_{i=1}^d\Gamma(\nu_i)a_i^{-\nu_i}}
      {\Gamma(\kappa+1)}.
 \label{eq:uniform-corner-constant}
\end{equation}
Equivalently, there are \(z_1>0\) and \(C<\infty\), depending
only on the constants in (G1)--(G3), such that
\[
 \sup_{\text{family}}
 r_z
 \left|
 \frac{
  \E[w\{\Phi(X)\}\mathbf 1_{\{D(\Phi(X))\le z\}}]
 }{
  w(0)C(\mu,\Sigma,a)
  z^\kappa(\sqrt{2\pi}r_z)^{\kappa-d}e^{-br_z}
 }-1
 \right|
 \le C
\]
for \(0<z\le z_1\).

In particular, if
\[
 H(z)
 =
 \E\left[
  w\{\Phi(X)\}\mathbf 1_{\{D(\Phi(X))\le z\}}
 \right],
\]
then there are uniform \(z_0,h_0,C_D>0\) such that
\[
 H(z_0)\ge h_0,
 \qquad
 H(2z)\le C_DH(z),
 \qquad 0<z\le z_0/2.
\]
Thus the weighted deficit masses satisfy condition {\rm(P3)}.
\end{theorem}

\begin{proof}
We divide the proof into five steps.

\medskip
\noindent
\emph{Step 1: density after corner rescaling.}
The density of \(U=\Phi(X)\) is
\begin{equation}
 f_U(u)
 =
 |\Sigma|^{-1/2}
 \exp\left\{
 -\frac12z(u)^\top(\Omega-I)z(u)
 +z(u)^\top\Omega\mu
 -\frac12\mu^\top\Omega\mu
 \right\},
 \label{eq:corner-U-density}
\end{equation}
where \(z_i(u)=\Phi^{-1}(u_i)\).

Put \(r=r_z\), and, for \(x_i>0\),
\[
 q_i=-\Phi^{-1}(zx_i).
\]
Thus \(z(zx)=-q\). Define
\begin{equation}
 A_z
 =
 |\Sigma|^{-1/2}
 e^{-\mu^\top\Omega\mu/2}
 z^\kappa
 (\sqrt{2\pi}r)^{\kappa-d}
 e^{-br},
 \label{eq:corner-normalizer}
\end{equation}
and
\[
 g_z(x)=\frac{z^df_U(zx)}{A_z}.
\]
By \eqref{eq:normal-log-tail},
\[
 \log z
 =
 -\frac{r^2}{2}
 -\log(\sqrt{2\pi}r)
 +\eta(r),
 \qquad |\eta(r)|\le r^{-2}.
\]
Substitution into \eqref{eq:corner-U-density} gives the exact
identity
\begin{align}
 \log g_z(x)
 &=
 -\frac12q^\top(\Omega-I)q
 -q^\top\Omega\mu
 +\frac12(\kappa-d)r^2
 +br
 +(d-\kappa)\eta(r).
 \label{eq:corner-log-g}
\end{align}

\medskip
\noindent
\emph{Step 2: a uniform integrable envelope near all coordinate
faces.}
Choose
\begin{equation}
 0<\vartheta<
 \frac12\min\{1,\nu_0,\omega_-\}.
 \label{eq:corner-theta}
\end{equation}
We prove that, for every fixed \(M<\infty\),
\begin{equation}
 g_z(x)\le C_M\prod_{i=1}^dx_i^{\vartheta-1},
 \qquad
 x\in(0,M]^d,
 \label{eq:corner-envelope}
\end{equation}
uniformly for small \(z\).

For each \(i\), \eqref{eq:q-implicit} gives
\[
 \log x_i
 =
 -\frac{q_i^2-r^2}{2}
 -\log\frac{q_i}{r}
 +\eta(q_i)-\eta(r).
\]
Since \(q_i\ge r-C_M/r\) by
Lemma~\ref{lem:uniform-normal-coordinate}, both
\(\eta(q_i)\) and \(\eta(r)\) are \(O(r^{-2})\), uniformly.
Combining this identity with \eqref{eq:corner-log-g} yields
\begin{align}
 &\log g_z(x)
 -(\vartheta-1)\sum_{i=1}^d\log x_i
 \nonumber\\
 &\quad=
 -\frac12\left[
 q^\top(\Omega-\vartheta I)q
 -r^2\mathbf 1^\top(\Omega-\vartheta I)\mathbf 1
 \right]
 -(q-r\mathbf 1)^\top\Omega\mu
 \nonumber\\
 &\qquad
 +(\vartheta-1)\sum_{i=1}^d\log\frac{q_i}{r}
 +O(r^{-2}).
 \label{eq:corner-envelope-log}
\end{align}
Write \(s=q-r\mathbf 1\) and \(B=\Omega-\vartheta I\).
By \eqref{eq:corner-theta},
\[
 \lambda_{\min}(B)\ge\omega_-/2,
 \qquad
 (B\mathbf 1)_i=\nu_i-\vartheta\ge\nu_0/2.
\]
Hence
\[
 -\frac12\{q^\top Bq-r^2\mathbf 1^\top B\mathbf 1\}
 =
 -r(B\mathbf 1)^\top s-\frac12s^\top Bs.
\]
The negative parts of the \(s_i\)'s are bounded by \(C/r\),
again by \eqref{eq:q-lower-uniform}. Therefore
\[
 -r(B\mathbf 1)^\top s
 \le
 -\frac{\nu_0r}{2}\sum_i s_i^+ +C.
\]
Also, by Young's inequality,
\[
 -\frac12s^\top Bs-s^\top\Omega\mu
 \le
 -\frac{\omega_-}{8}\|s\|_2^2+C.
\]
Finally, if \(q_i\ge r\), then
\((\vartheta-1)\log(q_i/r)\le0\), whereas if \(q_i<r\),
\[
 \left|\log\frac{q_i}{r}\right|\le Cr^{-2}.
\]
Thus the right side of \eqref{eq:corner-envelope-log} is bounded
above by a uniform constant. This proves
\eqref{eq:corner-envelope}.

In particular, on every fixed box,
\(\prod_i x_i^{\vartheta-1}\) is integrable. Moreover,
\begin{equation}
 \int_{\substack{x\in[0,M]^d\\
                 \min_i x_i<e^{-\sqrt r}}}
 \prod_{i=1}^dx_i^{\vartheta-1}\,dx
 \le Ce^{-\vartheta\sqrt r}
 =o(r^{-1}).
 \label{eq:corner-face-small}
\end{equation}
This proves uniform integrable domination at every coordinate
face, including intersections of several faces.

\medskip
\noindent
\emph{Step 3: the density expansion away from the thin face
region.}
Let
\[
 \mathcal E_r(M)
 =
 \left\{
  x\in(0,M]^d:\min_i x_i\ge e^{-\sqrt r}
 \right\},
 \qquad
 \ell_i=\log x_i.
\]
On this set,
\[
 -\sqrt r\le\ell_i\le\log M.
\]
Lemma~\ref{lem:uniform-normal-coordinate} gives
\[
 q
 =
 r\mathbf 1-\frac{\ell}{r}+e,
 \qquad
 |e_i|
 \le
 \frac{C_M(1+|\ell_i|+\ell_i^2)}{r^3}.
\]
Writing \(A=\Omega-I\) and expanding
\eqref{eq:corner-log-g}, one obtains
\begin{align}
 \log g_z(x)
 &=
 \sum_{i=1}^d(\nu_i-1)\ell_i
 +\frac{\ell^\top\Omega\mu}{r}
 -\frac{\ell^\top A\ell}{2r^2}
 +R_z(x),
 \label{eq:corner-density-expansion}
\end{align}
where
\begin{equation}
 |R_z(x)|
 \le
 \frac{C_M}{r^2}
 \left\{
  1+\sum_i|\ell_i|+\sum_i\ell_i^2
 \right\}.
 \label{eq:corner-density-R}
\end{equation}
Indeed, the terms containing \(e\) are bounded by
\[
 Cr\sum_i|e_i|
 +C\|e\|_2
 +C\frac{\|\ell\|_2\|e\|_2}{r}
 +C\|e\|_2^2,
\]
and \eqref{eq:corner-density-R} follows from \(d\le d_0\).

Consequently,
\begin{align}
 \left|
 \log\frac{g_z(x)}{\prod_ix_i^{\nu_i-1}}
 \right|
 &\le
 C_M\left[
 \frac{1+\sum_i|\log x_i|}{r}
 +\frac{1+\sum_i(\log x_i)^2}{r^2}
 \right].
 \label{eq:corner-log-relative}
\end{align}
The right side is \(O(r^{-1/2})\) on
\(\mathcal E_r(M)\). Hence, for large \(r\),
\begin{align}
 \left|
 g_z(x)-\prod_ix_i^{\nu_i-1}
 \right|
 &\le
 \frac{C_M}{r}
 \prod_ix_i^{\nu_i-1}
 \left\{
  1+\sum_i|\log x_i|
  +\sum_i(\log x_i)^2
 \right\}.
 \label{eq:corner-relative-error}
\end{align}
The parameters \(\nu_i\) range over a compact subset of
\((0,\infty)\): their lower bound is \(\nu_0\), while their
upper bound follows from \(d\le d_0\) and
\(\|\Omega\|_{\mathrm{op}}\le\omega_+\). Therefore
\[
 \sup_{\text{family}}
 \int_{[0,M]^d}
 \prod_ix_i^{\nu_i-1}
 \left\{
  1+\sum_i|\log x_i|
  +\sum_i(\log x_i)^2
 \right\}dx
 <\infty.
\]
Integrating \eqref{eq:corner-relative-error} and using
\eqref{eq:corner-face-small} gives
\begin{equation}
 \int_Bg_z(x)\,dx
 =
 \int_B\prod_ix_i^{\nu_i-1}\,dx+O(r^{-1})
 \label{eq:corner-integrated-density}
\end{equation}
uniformly for every measurable
\(B\subseteq[0,M]^d\). This is the required integrated
\(O(r^{-1})\) remainder.

\medskip
\noindent
\emph{Step 4: the moving nonlinear boundary.}
For sufficiently small \(z\), the event \(D(u)\le z\) is
contained in \([0,\eta_0]^d\), by the uniform gap \(d_*\).
On this event,
\[
 c_*\|u\|_1\le D(u)\le z,
\]
so after writing \(u=zx\),
\[
 \|x\|_1\le c_*^{-1}.
\]
Thus all rescaled domains lie in a common compact box.

Taylor's theorem and the Hessian bound give, uniformly on that
box,
\begin{equation}
 \frac{D(zx)}{z}
 =
 a^\top x+zR_D(z,x),
 \qquad
 |R_D(z,x)|\le C.
 \label{eq:corner-D-Taylor}
\end{equation}
Hence, for a uniform \(C_0<\infty\),
\begin{equation}
 \{x\ge0:a^\top x\le1-C_0z\}
 \subseteq
 \left\{x\ge0:\frac{D(zx)}z\le1\right\}
 \subseteq
 \{x\ge0:a^\top x\le1+C_0z\}.
 \label{eq:corner-domain-sandwich}
\end{equation}

Put
\[
 J(t)
 =
 \int_{\{x\ge0:a^\top x\le t\}}
 \prod_ix_i^{\nu_i-1}\,dx.
\]
The change of variables \(x=ty\) gives
\[
 J(t)=t^\kappa J(1),
\]
where
\begin{equation}
 J(1)
 =
 \frac{\prod_i\Gamma(\nu_i)a_i^{-\nu_i}}
      {\Gamma(\kappa+1)}.
 \label{eq:corner-dirichlet-integral}
\end{equation}
The parameters \(\nu_i,a_i,\kappa\) range over compact sets, so
\(J(1)\) is uniformly bounded above and away from zero.
The sandwich \eqref{eq:corner-domain-sandwich} therefore gives
\begin{equation}
 \int_{\{D(zx)\le z\}}\prod_ix_i^{\nu_i-1}\,dx
 =
 J(1)+O(z)
 \label{eq:corner-boundary-error}
\end{equation}
uniformly. This proves stability under the moving nonlinear
boundary, including its intersections with the coordinate
faces.

Combining \eqref{eq:corner-integrated-density} and
\eqref{eq:corner-boundary-error},
\begin{equation}
 \int_{\{D(zx)\le z\}}g_z(x)\,dx
 =
 J(1)+O(r^{-1}),
 \label{eq:corner-unweighted-integral}
\end{equation}
because \(z=o(r^{-1})\).

\medskip
\noindent
\emph{Step 5: the weight and the final constant.}
On the rescaled event,
\[
 \|zx\|_1\le z/c_*,
\]
so (G3) gives
\[
 w(zx)=w(0)+O(z)
\]
uniformly. The envelope \eqref{eq:corner-envelope} is
integrable, and \(w(0)\) is bounded above and away from zero.
Thus
\begin{align}
 \int_{\{D(zx)\le z\}}w(zx)g_z(x)\,dx
 &=
 w(0)J(1)+O(r^{-1}).
 \label{eq:corner-weighted-integral}
\end{align}
Multiplication by the normalizer \(A_z\) in
\eqref{eq:corner-normalizer}, followed by
\eqref{eq:corner-dirichlet-integral}, proves
\eqref{eq:uniform-corner-main} and
\eqref{eq:uniform-corner-constant}.

It remains to prove uniform doubling. By
Lemma~\ref{lem:uniform-normal-coordinate}, applied with fixed
\(x=2\),
\begin{equation}
 r_{2z}
 =
 r_z-\frac{\log2}{r_z}+O(r_z^{-3}).
 \label{eq:r2z-expansion}
\end{equation}
Therefore
\[
 \left(\frac{r_{2z}}{r_z}\right)^{\kappa-d}
 \exp\{-b(r_{2z}-r_z)\}
\]
is uniformly bounded above and away from zero for small \(z\).
Since \(\kappa\) is uniformly bounded,
\[
 \frac{H(2z)}{H(z)}
 =
 2^\kappa
 \left(\frac{r_{2z}}{r_z}\right)^{\kappa-d}
 e^{-b(r_{2z}-r_z)}
 \{1+O(r_z^{-1})\}
\]
is uniformly bounded. Choosing one sufficiently small \(z_0\)
also makes the relative error in
\eqref{eq:uniform-corner-main} smaller than \(1/2\). The leading
constant at this fixed \(z_0\) has a positive uniform lower
bound, proving \(H(z_0)\ge h_0\).
\end{proof}


\subsection{Uniform verification for the conditional Gaussian families}

Let
\[
 R_4=(\rho^{|k-\ell|})_{1\le k,\ell\le4},
 \qquad \rho=\frac14.
\]
Its precision matrix is
\begin{equation}
 Q=R_4^{-1}
 =
 \frac1{1-\rho^2}
 \begin{pmatrix}
 1&-\rho&0&0\\
 -\rho&1+\rho^2&-\rho&0\\
 0&-\rho&1+\rho^2&-\rho\\
 0&0&-\rho&1
 \end{pmatrix},
 \label{eq:R4-precision}
\end{equation}
and
\begin{equation}
 Q\mathbf 1
 =
 \left(\frac45,\frac35,\frac35,\frac45\right)^\top.
 \label{eq:R4-row-sums}
\end{equation}

\begin{corollary}[The corner theorem applies uniformly to A1--A3 and B1]
\label{cor:corner-model-families}
Fix \(0<\varepsilon<1/2\). For every coordinate \(j\), every
conditioning value
\[
 u\in I_{\varepsilon/2},
\]
and every model A1--A3 or B1, the conditional Gaussian family of
the free active latent coordinates satisfies
Theorem~\ref{thm:uniform-gaussian-corner}, with
\[
 \nu_0=\frac35.
\]
Its exponents are
\begin{equation}
 \kappa_j=
 \begin{cases}
  34/15,&j=1,4,\\
  41/15,&j=2,3,\\
  14/5+\lambda_j^2/(1-\lambda_j^2),
       &j>4\text{ in the Gaussian AR block},\\
  14/5,&j\in\{5,\ldots,24\}\text{ in B1},
 \end{cases}
 \label{eq:corner-kappas-rigorous}
\end{equation}
where
\[
 \lambda_j=\rho^{j-4}.
\]
In particular,
\[
 \frac{34}{15}\le\kappa_j\le\frac{43}{15}.
\]
The B1 deficit mass is a weighted Gaussian-corner mass whose
weight satisfies (G3) uniformly.
\end{corollary}

\begin{proof}
We first identify all conditional precision matrices.

If \(j\in\{1,2,3,4\}\), the conditional precision matrix of
\(Z_{\{1,2,3,4\}\setminus\{j\}}\) given \(Z_j\) is the principal
submatrix \(Q_{-j,-j}\). For \(j=1\), its row-sum vector is
\[
 \left(\frac{13}{15},\frac35,\frac45\right)^\top,
\]
whose sum is \(34/15\). For \(j=2\), it is
\[
 \left(\frac{16}{15},\frac{13}{15},\frac45\right)^\top,
\]
whose sum is \(41/15\). The cases \(j=3,4\) are symmetric.
Every component is at least \(3/5\).

Now let \(j>4\) belong to the original Gaussian AR block.
Writing \(\lambda=\rho^{j-4}\), the covariance between
\((Z_1,\ldots,Z_4)\) and \(Z_j\) is
\[
 \lambda R_4e_4.
\]
The conditional covariance of the active block is therefore
\[
 \Sigma_\lambda
 =
 R_4-\lambda^2R_4e_4e_4^\top R_4.
\]
The Sherman--Morrison identity gives
\begin{equation}
 \Sigma_\lambda^{-1}
 =
 Q+\frac{\lambda^2}{1-\lambda^2}e_4e_4^\top.
 \label{eq:inactive-conditional-precision}
\end{equation}
Hence
\[
 \Sigma_\lambda^{-1}\mathbf 1
 =
 \left(
  \frac45,\frac35,\frac35,
  \frac45+\frac{\lambda^2}{1-\lambda^2}
 \right)^\top,
\]
and
\[
 \kappa
 =
 \frac{14}{5}+\frac{\lambda^2}{1-\lambda^2}.
\]
Since \(0\le\lambda\le\rho=1/4\),
\[
 \frac{\lambda^2}{1-\lambda^2}\le\frac1{15},
\]
which gives the upper bound \(43/15\).

In B1, a proxy coordinate is independent of the active Gaussian
block. Conditioning on that proxy leaves covariance \(R_4\),
so the precision is \(Q\) and \(\kappa=14/5\). A background
coordinate \(j\ge25\) gives the family
\eqref{eq:inactive-conditional-precision}, with the still
smaller range \(\lambda\le\rho^{21}\).

The conditional means are also uniform. If \(j\le4\), the mean
is a fixed covariance vector times
\(\Phi^{-1}(u)\). If \(j>4\) is in the AR block, it is
\[
 \lambda R_4e_4\Phi^{-1}(u).
\]
Since \(u\in I_{\varepsilon/2}\),
\(\Phi^{-1}(u)\) ranges over a bounded interval. The covariance
families just displayed are compact and uniformly positive
definite. Thus (G1) holds.

It remains to verify the deficit functions. Write
\[
 \frac1{\gamma(x)}=2e^{\psi(x_1,\ldots,x_4)},
\]
where
\[
 \psi_{A1}(x)=\psi_{B1}(x)=\sum_{\ell=1}^4x_\ell,
\]
\[
 \psi_{A2}(x)=\frac12\sum_{\ell=1}^4x_\ell,
\]
and
\[
 \psi_{A3}(x)
 =
 0.8\sum_{\ell=1}^4x_\ell
 +0.20x_1x_2+0.15x_3x_4.
\]
Fixing the conditioning coordinate and denoting the free active
coordinates by \(y\), let
\[
 \psi_0=\psi\big|_{y=0}.
\]
Then the reciprocal-index deficit is
\begin{equation}
 D(y)
 =
 2e^{\psi_0}
 \left[
  e^{\psi(y)-\psi_0}-1
 \right].
 \label{eq:model-deficit-form}
\end{equation}
For A1 and B1, every free-coordinate derivative of \(\psi\) is
one. For A2 it is \(1/2\). For A3, every such derivative at the
free corner lies in \([0.8,1]\): a derivative is either \(0.8\),
\(0.8+0.20u\), or \(0.8+0.15u\). Thus every component of
\(\nabla D(0)\) is uniformly bounded above and away from zero.
The Hessians are uniformly bounded, and strict coordinatewise
monotonicity shows that \(D(y)=0\) only at \(y=0\). Compactness
of the remaining conditioning parameters gives the uniform
two-sided corner comparison and the uniform gap outside a fixed
corner neighbourhood. Hence (G2) holds.

For A1--A3 the weight is one. For B1, let
\[
 A(F)=\frac{\Phi(F)-1/2}{2},
 \qquad |A(F)|\le\frac14.
\]
The leading conditional tail coefficient is
\[
 c_B(x)=
 \E\left[
  e^{A(F)/\gamma(x)}
  \,\middle|\,U=x
 \right].
\]
For the deficit mass conditional only on \(U_j=u\), the tower
property gives
\begin{align}
 &\E\left[
  c_B(U)\mathbf 1_{\{D(U_{\mathcal A})\le z\}}
  \,\middle|\,U_j=u
 \right]
 \nonumber\\
 &\qquad=
 \E\left[
  e^{A(F)/\gamma(U)}
  \mathbf 1_{\{D(U_{\mathcal A})\le z\}}
  \,\middle|\,U_j=u
 \right].
 \label{eq:B1-weight-tower}
\end{align}
After conditioning on the free active vector \(y\), this is a
Gaussian-corner integral with weight
\[
 w_{j,u}(y)
 =
 \E\left[
  e^{h(y)A(F)}
  \,\middle|\,U_j=u,\;U_{\mathcal A\setminus\{j\}}=y
 \right],
 \qquad
 h(y)=\frac1{\gamma(y)}.
\]
If \(j\) is not a proxy, \(F\) is standard normal and independent
of the conditioning variables. If \(j\) is a proxy and
\(t=\Phi^{-1}(u)\), then
\[
 F\mid Z_j=t\sim N(0.7t,\,1-0.7^2).
\]
These means are uniformly bounded on \(I_{\varepsilon/2}\).

For any Gaussian law of \(F\), define
\[
 W(h)=\E e^{hA(F)}.
\]
Differentiation under the expectation gives
\[
 \left|\frac{d}{dh}\log W(h)\right|
 =
 \left|
 \frac{\E[A(F)e^{hA(F)}]}{\E e^{hA(F)}}
 \right|
 \le\frac14.
\]
In B1,
\[
 2\le h(y)\le2e^4,
\]
and \(h\) has a uniform bounded derivative in the active PIT
coordinates. Therefore
\[
 |\log w_{j,u}(y)-\log w_{j,u}(0)|
 \le C\|y\|_1.
\]
Also
\[
 e^{-e^4/2}\le w_{j,u}(0)\le e^{e^4/2}.
\]
It follows that
\[
 |w_{j,u}(y)-w_{j,u}(0)|\le C\|y\|_1
\]
uniformly in \(j,u\). This is (G3), and the weighted conclusion
of Theorem~\ref{thm:uniform-gaussian-corner} applies.
\end{proof}


\subsection{Exact primitive tail representation for B1}

We now prove the horizontal-coupling assertion with explicit
constants.

Let
\[
 \lambda=0.7,\qquad
 \sigma^2=1-\lambda^2=0.51,
 \qquad
 \mathcal P=\{5,\ldots,24\},
 \qquad q=|\mathcal P|=20.
\]
The proxy coordinates are
\[
 Z_\ell=\lambda F+\sigma\epsilon_\ell,
 \qquad \ell\in\mathcal P,
\]
where \(F,\epsilon_5,\ldots,\epsilon_{24}\) are independent
standard normals. The proxy block is independent of the base
Gaussian AR block.

The B1 response is generated by
\begin{equation}
 Y
 =
 S^{\gamma(U)}
 \exp\left\{
  A(F)-\frac1{2S}
 \right\},
 \qquad
 S=V^{-1},\quad V\sim{\rm Unif}(0,1),
 \label{eq:B1-generator-rigorous}
\end{equation}
where
\[
 A(f)=\frac{\Phi(f)-1/2}{2}.
\]
Put
\[
 h(x)=\frac1{\gamma(x)}
 =
 2\exp\left(\sum_{a=1}^4x_a\right),
 \qquad
 2\le h(x)\le h_+:=2e^4.
\]

Conditioning on the complete observed Gaussian vector
\(Z=z\), the posterior law of \(F\) is
\begin{equation}
 F\mid Z=z
 \sim N\{m(z),s_*^2\},
 \label{eq:B1-posterior}
\end{equation}
where
\begin{equation}
 s_*^2
 =
 \frac{\sigma^2}{\sigma^2+q\lambda^2},
 \qquad
 m(z)
 =
 \beta\sum_{\ell\in\mathcal P}z_\ell,
 \qquad
 \beta=
 \frac{\lambda}{\sigma^2+q\lambda^2}.
 \label{eq:B1-posterior-parameters}
\end{equation}

Let \(W_0\) denote the principal Lambert \(W\)-function and set
\[
 R(x)=
 \begin{cases}
  W_0(x)/x,&x>0,\\
  1,&x=0.
 \end{cases}
\]
Since \(x=W_0(x)e^{W_0(x)}\),
\begin{equation}
 R(x)=e^{-W_0(x)}.
 \label{eq:Lambert-ratio}
\end{equation}

For \(h\in[2,h_+]\), \(a\in[-1/4,1/4]\), and \(t\ge1\), define
\begin{equation}
 x_t(h,a)
 =
 \frac h2e^{-h(t-a)}
 \label{eq:B1-x}
\end{equation}
and
\begin{equation}
 G_t(h,a)
 =
 e^{ha}R\{x_t(h,a)\}
 =
 \exp\left[
  ha-W_0\{x_t(h,a)\}
 \right].
 \label{eq:B1-G}
\end{equation}
Finally, for \(G\sim N(0,1)\), define
\begin{align}
 c(h,m)
 &=
 \E\exp\{hA(m+s_*G)\},
 \label{eq:B1-c-hm}\\
 H_t(h,m)
 &=
 \E G_t\{h,A(m+s_*G)\}.
 \label{eq:B1-H-hm}
\end{align}

\begin{lemma}[Exact B1 tail factorization]
\label{lem:B1-exact-tail}
For every observed \(x\in(0,1)^p\), let
\(z=\Phi^{-1}(x)\) coordinatewise and put
\[
 c_B(x)=c\{h(x),m(z)\},
 \qquad
 1+r_B(t,x)
 =
 \frac{H_t\{h(x),m(z)\}}{c\{h(x),m(z)\}}.
\]
Then, for every \(t\ge1\),
\begin{equation}
 \Pp(Y>e^t\mid U=x)
 =
 c_B(x)e^{-t/\gamma(x)}
 \{1+r_B(t,x)\}
 \label{eq:B1-exact-factorization}
\end{equation}
exactly. Moreover,
\begin{equation}
 -e^{1/2-2t}\le r_B(t,x)\le0,
 \label{eq:B1-r-bound}
\end{equation}
and, wherever the derivative is taken,
\begin{equation}
 |\partial_t r_B(t,x)|
 \le h_+e^{1/2-2t}.
 \label{eq:B1-r-derivative}
\end{equation}
Thus B1 satisfies (P2), with a remainder and log-level derivative
of order \(O(e^{-2t})\), uniformly in \(p\) and \(x\).
\end{lemma}

\begin{proof}
Fix \(x,F\), write \(g=\gamma(x)\), \(h=1/g\),
\(a=A(F)\), and let \(p\) be the conditional survival
probability at \(e^t\). By monotonicity of the generator,
\(p\) is the unique solution in \((0,1)\) of
\[
 e^t=p^{-g}e^{a-p/2}.
\]
Equivalently,
\[
 p\,e^{hp/2}=e^{-h(t-a)}.
\]
Multiplying by \(h/2\) and applying \(W_0\),
\[
 p
 =
 \frac2hW_0\left\{\frac h2e^{-h(t-a)}\right\}.
\]
Using \eqref{eq:Lambert-ratio},
\begin{align}
 p
 &=
 e^{-h(t-a)}
 R\left\{\frac h2e^{-h(t-a)}\right\}
 \nonumber\\
 &=
 e^{-ht}G_t(h,a).
 \label{eq:B1-p-exact}
\end{align}
Averaging \eqref{eq:B1-p-exact} under the posterior
\eqref{eq:B1-posterior} gives
\[
 \Pp(Y>e^t\mid U=x)
 =
 e^{-h(x)t}H_t\{h(x),m(z)\},
\]
which is \eqref{eq:B1-exact-factorization}.

For \(t\ge1\), \(a\le1/4\), and \(h\ge2\), put
\(y=h(t-a)\). Then \(y\ge3/2\), and the function
\(h\mapsto(h/2)e^{-h(t-a)}\) is decreasing on
\([2,\infty)\). Hence
\begin{equation}
 0\le x_t(h,a)\le e^{-2(t-1/4)}
 =e^{1/2-2t}.
 \label{eq:B1-x-bound}
\end{equation}
Since \(0\le W_0(x)\le x\),
\[
 1-x\le e^{-W_0(x)}\le1.
\]
The quotient \(H_t/c\) is a weighted average of
\(R(x_t)=e^{-W_0(x_t)}\), with weights proportional to
\(e^{hA(F)}\). Therefore
\[
 1-e^{1/2-2t}
 \le \frac{H_t(h,m)}{c(h,m)}\le1,
\]
which proves \eqref{eq:B1-r-bound}.

Finally,
\[
 \partial_t\log R\{x_t(h,a)\}
 =
 \frac{hW_0\{x_t(h,a)\}}
      {1+W_0\{x_t(h,a)\}},
\]
so
\[
 0\le\partial_tR\{x_t(h,a)\}
 \le h\,x_t(h,a).
\]
Averaging with the same fixed weights gives
\[
 0\le\partial_t r_B(t,x)
 \le h_+\sup_{h,a}x_t(h,a)
 \le h_+e^{1/2-2t}.
\]
\end{proof}


\subsection{Uniform horizontal coupling for B1}

Let
\[
 T_\varepsilon=\Phi^{-1}(1-\varepsilon/2),
 \qquad
 L_\varepsilon=\frac1{\varphi(T_\varepsilon)},
 \qquad
 \varphi_0=\varphi(0)=\frac1{\sqrt{2\pi}}.
\]
Also put
\[
 a_0=\frac14,
 \qquad
 a_1=\sup_f|A'(f)|
 =\frac{\varphi_0}{2},
\]
and
\begin{equation}
 C_h
 =
 \frac14+\frac54e^{-3/2}.
 \label{eq:B1-Ch}
\end{equation}

\begin{proposition}[Uniform B1 horizontal coupling]
\label{prop:B1-horizontal-coupling}
For every \(p\ge25\), every coordinate \(j\le p\),
every
\[
 u\in I_\varepsilon,
 \qquad
 v\in I_{\varepsilon/2},
 \qquad
 |u-v|\le\varepsilon/2,
\]
the conditional kernels \(K_j(u,\cdot)\) and
\(K_j(v,\cdot)\) admit a coupling
\[
 X_u\sim\mathcal L(U\mid U_j=u),
 \qquad
 X_v\sim\mathcal L(U\mid U_j=v)
\]
such that, almost surely,
\begin{equation}
 \|X_v-X_u\|_\infty
 \le
 \varphi_0L_\varepsilon|v-u|,
 \label{eq:B1-U-coupling}
\end{equation}
and
\begin{equation}
 \left|
 \frac1{\gamma(X_v)}-\frac1{\gamma(X_u)}
 \right|
 \le L_D|v-u|,
 \label{eq:B1-LD}
\end{equation}
where
\begin{equation}
 L_D=8e^4\varphi_0L_\varepsilon.
 \label{eq:B1-LD-constant}
\end{equation}

For the exact factorization
\eqref{eq:B1-exact-factorization}, one also has, almost surely
and for every \(t\ge1\),
\begin{align}
 &|\log c_B(X_v)-\log c_B(X_u)|
 \nonumber\\
 &\quad+
 |\log\{1+r_B(t,X_v)\}
   -\log\{1+r_B(t,X_u)\}|
 \le L_H|v-u|,
 \label{eq:B1-LH}
\end{align}
where the deterministic constant
\begin{equation}
 L_H
 =
 \left(
  \frac34+\frac54e^{-3/2}
 \right)L_D
 +3e^4\varphi_0\lambda L_\varepsilon
 \label{eq:B1-LH-constant}
\end{equation}
is independent of \(p,n,j,u,v,t\).

Thus B1 satisfies condition {\rm(P4)} with
\[
 t_0=1,\qquad \zeta_n=0.
\]
\end{proposition}

\begin{proof}
We first construct the coupling, then prove deterministic
Lipschitz bounds for the two tail factors.

\medskip
\noindent
\emph{Step 1: common-residual Gaussian coupling.}
Let \(R_p\) denote the correlation matrix of the observed
Gaussian vector \(Z=(Z_1,\ldots,Z_p)\). It has two independent
blocks:

\begin{itemize}
\item on
\[
 \mathcal B=\{1,2,3,4,25,\ldots,p\},
\]
\[
 (R_p)_{k\ell}=\rho^{|k-\ell|};
\]

\item on the proxy block \(\mathcal P\),
\[
 (R_p)_{kk}=1,
 \qquad
 (R_p)_{k\ell}=\lambda^2
 \quad(k\ne\ell);
\]

\item correlations between \(\mathcal B\) and \(\mathcal P\)
are zero.
\end{itemize}

Put
\[
 t_u=\Phi^{-1}(u),
 \qquad
 t_v=\Phi^{-1}(v).
\]
Let \(r_j=(R_p)_{-j,j}\), and let
\[
 \eta_j\sim
 N\left(
  0,\,
  (R_p)_{-j,-j}-r_jr_j^\top
 \right).
\]
Using the same \(\eta_j\), define
\[
 Z^{(u)}_j=t_u,\qquad
 Z^{(u)}_{-j}=r_jt_u+\eta_j,
\]
\[
 Z^{(v)}_j=t_v,\qquad
 Z^{(v)}_{-j}=r_jt_v+\eta_j.
\]
These are the correct conditional Gaussian laws given
\(Z_j=t_u\) and \(Z_j=t_v\). Coordinatewise,
\begin{equation}
 Z^{(v)}_k-Z^{(u)}_k
 =
 (R_p)_{kj}(t_v-t_u).
 \label{eq:B1-Z-difference}
\end{equation}
Every correlation has absolute value at most one. Since
\(u,v\in I_{\varepsilon/2}\), the inverse-function theorem gives
\begin{equation}
 |t_v-t_u|
 \le L_\varepsilon|v-u|.
 \label{eq:B1-t-difference}
\end{equation}
Consequently,
\[
 \|Z^{(v)}-Z^{(u)}\|_\infty
 \le L_\varepsilon|v-u|.
\]
Applying the global Lipschitz bound
\(|\Phi(s)-\Phi(t)|\le\varphi_0|s-t|\) proves
\eqref{eq:B1-U-coupling} for
\[
 X_u=\Phi\{Z^{(u)}\},
 \qquad
 X_v=\Phi\{Z^{(v)}\}.
\]

\medskip
\noindent
\emph{Step 2: reciprocal-index constant.}
Writing \(z_a\) for the first four latent coordinates,
\[
 h(z)=\frac1{\gamma\{\Phi(z)\}}
 =
 2\exp\left\{\sum_{a=1}^4\Phi(z_a)\right\}.
\]
For \(a\le4\),
\[
 \left|\frac{\partial h}{\partial z_a}\right|
 =
 h(z)\varphi(z_a)
 \le2e^4\varphi_0.
\]
Therefore
\begin{align}
 |h\{Z^{(v)}\}-h\{Z^{(u)}\}|
 &\le
 8e^4\varphi_0
 \|Z_{\{1,\ldots,4\}}^{(v)}
   -Z_{\{1,\ldots,4\}}^{(u)}\|_\infty
 \nonumber\\
 &\le
 8e^4\varphi_0L_\varepsilon|v-u|.
\end{align}
This is \eqref{eq:B1-LD}--\eqref{eq:B1-LD-constant}.

\medskip
\noindent
\emph{Step 3: movement of the posterior mean of \(F\).}
Recall
\[
 m(z)=\beta\sum_{\ell\in\mathcal P}z_\ell,
 \qquad
 \beta=\frac{\lambda}{\sigma^2+q\lambda^2}.
\]
Under the common-residual coupling,
\[
 m\{Z^{(v)}\}-m\{Z^{(u)}\}
 =
 \beta(t_v-t_u)
 \sum_{\ell\in\mathcal P}(R_p)_{\ell j}.
\]
If \(j\notin\mathcal P\), the sum is zero. If
\(j\in\mathcal P\), then
\[
 \sum_{\ell\in\mathcal P}(R_p)_{\ell j}
 =
 1+(q-1)\lambda^2
 =
 \sigma^2+q\lambda^2.
\]
Hence, in all cases,
\begin{equation}
 |m\{Z^{(v)}\}-m\{Z^{(u)}\}|
 \le
 \lambda|t_v-t_u|
 \le
 \lambda L_\varepsilon|v-u|.
 \label{eq:B1-m-difference}
\end{equation}
This is where the fact that the proxy-block size \(q=20\) is
fixed enters. In fact its contribution cancels exactly in the
posterior coefficient.

\medskip
\noindent
\emph{Step 4: derivatives of \(c(h,m)\).}
Write \(F_m=m+s_*G\), with \(G\sim N(0,1)\). Since
\(|A|\le a_0\) and \(|A'|\le a_1\), differentiation under the
expectation in \eqref{eq:B1-c-hm} gives
\begin{equation}
 \partial_h\log c(h,m)
 =
 \frac{\E[A(F_m)e^{hA(F_m)}]}
      {\E e^{hA(F_m)}},
 \qquad
 |\partial_h\log c(h,m)|\le a_0,
 \label{eq:B1-c-h-derivative}
\end{equation}
and
\begin{equation}
 \partial_m\log c(h,m)
 =
 \frac{\E[hA'(F_m)e^{hA(F_m)}]}
      {\E e^{hA(F_m)}},
 \qquad
 |\partial_m\log c(h,m)|\le h_+a_1.
 \label{eq:B1-c-m-derivative}
\end{equation}
Thus, for any two parameter pairs,
\begin{equation}
 |\log c(h_2,m_2)-\log c(h_1,m_1)|
 \le
 a_0|h_2-h_1|+h_+a_1|m_2-m_1|.
 \label{eq:B1-c-Lipschitz}
\end{equation}

\medskip
\noindent
\emph{Step 5: derivatives of \(H_t(h,m)\).}
For \(G_t\) in \eqref{eq:B1-G}, put
\[
 x=x_t(h,a),\qquad W=W_0(x).
\]
The identities
\[
 xW_0'(x)=\frac{W}{1+W},
 \qquad
 \partial_h\log x=\frac1h-(t-a),
 \qquad
 \partial_a\log x=h
\]
give
\begin{equation}
 \partial_h\log G_t(h,a)
 =
 a-\frac{W}{1+W}
 \left\{\frac1h-(t-a)\right\},
 \label{eq:B1-G-h}
\end{equation}
and
\begin{equation}
 \partial_a\log G_t(h,a)
 =
 \frac{h}{1+W}.
 \label{eq:B1-G-a}
\end{equation}
As in the proof of Lemma~\ref{lem:B1-exact-tail}, let
\(y=h(t-a)\ge3/2\). Since \(W\le x=(h/2)e^{-y}\),
\begin{align}
 |\partial_h\log G_t(h,a)|
 &\le
 a_0+\frac{x}{h}+x(t-a)
 \nonumber\\
 &=
 a_0+\frac12e^{-y}+\frac y2e^{-y}
 \nonumber\\
 &\le
 \frac14+\frac54e^{-3/2}
 =C_h.
 \label{eq:B1-G-h-bound}
\end{align}
Also
\begin{equation}
 |\partial_a\log G_t(h,a)|\le h_+.
 \label{eq:B1-G-a-bound}
\end{equation}

Because \(H_t(h,m)\) is the expectation of the positive
function \(G_t\{h,A(F_m)\}\), its logarithmic derivatives are
weighted averages of the corresponding logarithmic derivatives
of the integrand. Thus
\begin{equation}
 |\partial_h\log H_t(h,m)|\le C_h
 \label{eq:B1-H-h-derivative}
\end{equation}
and
\begin{equation}
 |\partial_m\log H_t(h,m)|
 \le h_+a_1.
 \label{eq:B1-H-m-derivative}
\end{equation}
The bounds hold uniformly for every \(t\ge1\).

Since
\[
 \log\{1+r_B(t,x)\}
 =
 \log H_t\{h(x),m(x)\}
 -\log c\{h(x),m(x)\},
\]
we obtain
\begin{align}
 &|\log\{1+r_B(t,x_2)\}
   -\log\{1+r_B(t,x_1)\}|
 \nonumber\\
 &\qquad\le
 (C_h+a_0)|h_2-h_1|
 +2h_+a_1|m_2-m_1|.
 \label{eq:B1-r-Lipschitz}
\end{align}
Adding \eqref{eq:B1-c-Lipschitz} gives
\begin{align}
 &|\log c_2-\log c_1|
 +|\log(1+r_2)-\log(1+r_1)|
 \nonumber\\
 &\qquad\le
 (C_h+2a_0)|h_2-h_1|
 +3h_+a_1|m_2-m_1|.
 \label{eq:B1-combined-Lipschitz}
\end{align}
Now use
\[
 C_h+2a_0
 =
 \frac34+\frac54e^{-3/2},
 \qquad
 3h_+a_1
 =
 3e^4\varphi_0,
\]
together with \eqref{eq:B1-LD} and
\eqref{eq:B1-m-difference}. This gives
\[
 |\log c_B(X_v)-\log c_B(X_u)|
 +
 |\log\{1+r_B(t,X_v)\}
  -\log\{1+r_B(t,X_u)\}|
 \le L_H|v-u|
\]
with \(L_H\) exactly as in
\eqref{eq:B1-LH-constant}. No exceptional deterministic term is
needed, so \(\zeta_n=0\).
\end{proof}


\begin{remark}[What has been repaired]
\label{rem:corner-coupling-repairs}
The preceding arguments make six points explicit.

\begin{enumerate}
\item The Gaussian density is dominated uniformly all the way to
intersections of coordinate faces by
\(\prod_i x_i^{\vartheta-1}\).

\item The normal-quantile expansion is used only off a face region
whose normalized mass is \(o(r_z^{-1})\).

\item The pointwise density error need not be \(O(r_z^{-1})\):
the mean term is proportional to
\(r_z^{-1}\sum_i|\log x_i|\). Its integral is
\(O(r_z^{-1})\), which is the statement actually required.

\item The nonlinear boundary contributes only \(O(z)\), because it
is trapped between two homothetic simplices. Since
\(z=o(r_z^{-1})\), it is negligible at the announced order.

\item A weight \(w\) changes the leading constant to
\(w(0)C(\mu,\Sigma,a)\). Omitting \(w(0)\) would be incorrect.

\item For B1, the horizontal constants are independent of the
ambient dimension. The common-residual coupling moves every
latent Gaussian coordinate by at most
\(L_\varepsilon|u-v|\); the only apparent sum over the twenty
proxies cancels exactly in the posterior mean coefficient.
\end{enumerate}
\end{remark}